\documentclass[10pt,a4paper]{amsart}
\usepackage[margin=19mm]{geometry}
\usepackage{amsmath,amssymb,mathtools}
\usepackage[scr=rsfs,cal=euler]{mathalfa}
\usepackage{xfrac}
\usepackage[unicode,hidelinks,bookmarks=false]{hyperref}
\newcommand{\ie}{i.e.\ }
\newcommand{\cf}{cf.\ }
\newcommand{\resp}{respectively\ }
\newcommand{\Subsection}{\subsection}
\providecommand{\nobreakdash}{\nobreak\hbox{-}}
\setlength{\emergencystretch}{2em}
\multlinegap0pt
\usepackage{bm}
\usepackage{bbm}
\usepackage{dsfont}
\usepackage{xspace}
\usepackage{cancel}
\usepackage{mathtools}
\usepackage{amsthm}
\makeatletter
\@ifundefined{theorem}{\newtheorem{theorem}{Theorem}}{}
\@ifundefined{lemma}{\newtheorem{lemma}[theorem]{Lemma}}{}
\makeatother
\makeatletter
\newcommand{\CC}{\mathbb{C}}
\newcommand{\Ind}{\mathrm{Ind}}
\newcommand{\ZZ}{\mathbb{Z}}
\expandafter\ifx\csname psmatrix\endcsname\relax
\newenvironment{psmatrix}
  {\left(\begin{smallmatrix}}
\fi
\expandafter\ifx\csname remarks\endcsname\relax
\newenvironment{remarks}
{\medskip\noindent{\it
Remarks.}\begin{list}{{\rm(\arabic{remarkscounter})}
}{\usecounter{remarkscounter}
\renewcommand{\theremarkscounter}{{\bf(\alph{remarkscounter})}}
\setlength{\labelsep}{\fill} \setlength{\leftmargin}{0pt}
\setlength{\itemindent}{\fill}
\setlength{\labelwidth}{\fill}\setlength{\topsep}{0pt}
\setlength{\listparindent}{0pt}}} {\end{list}}
\fi
\makeatother
\title[Schwartz spaces on L-monoids: non-Archimedean]{Schwartz spaces on L-monoids: non-Archimedean}
\author{Chun-Hsien Hsu, HaoYun Yao}
\date{Source: arXiv:2607.28507v2 (CC BY 4.0)}
\begin{document}
\maketitle

\noindent\emph{Source and licence.} Schwartz spaces on L-monoids: non-Archimedean. Version arXiv:2607.28507v2 (CC BY 4.0), distributed under the Creative Commons Attribution 4.0 International licence, as recorded in the arXiv licence metadata for this exact version and shown on the abstract page https://arxiv.org/abs/2607.28507v2. Full rights notice: licenses/arxiv-2607.28507-CC-BY-4.0.txt.\par\smallskip

\noindent\textit{Editorial scope.} Counted unit: the Lemma whose source label is \texttt{lem:varying} in the article (source0.tex lines 2692--2708, source label \texttt{lem:varying}), delivered together with the article's own complete original proof of it (source0.tex lines 2710--2788, one \texttt{proof} environment of 79 lines). No mathematics is written, repaired or re-derived: statement and proof are the article's own text line for line. Required context retained uncounted: lem:cuspnonzero (lines 2664--2669). Every definition, display and label of those lines is retained unchanged. Omitted from this package and not redistributed: 1--2663, 2670--2691, 2709--2709, 2789--3965. Nothing inside a retained excerpt is omitted or altered: every retained body is delivered exactly as the article's lines are written, so no omission receipt of any kind accompanies this package. Citations: the retained excerpts carry no citation command at all, so no bibliography entry of the article is transcribed and no reference is redistributed. Reference adaptation: none was needed and none was made; every reference inside the delivered unit resolves inside it, so no unresolved reference or citation is produced. Printed slips kept literal: no printed word, symbol, hypothesis or number of the counted statement or of its proof was altered. Layout and class: the article's own class is \texttt{amsart} with packages including amsmath, amsthm, amscd, amsfonts, amssymb, epic, eepic, bbm, comment, enumitem, mathdots, mathtools, mathrsfs, hyperref; the wrapper is the lane's pinned \texttt{amsart} editorial wrapper at 10pt on A4, which restates the theorem-like environments the retained text needs and adds the article's own macro definitions that the retained text needs (\texttt{\textbackslash CC}, \texttt{\textbackslash Ind}, \texttt{\textbackslash ZZ}), with extra wrapper packages (bm, bbm, dsfont, xspace, cancel, mathtools). The delivered numbering is the unit's own. Assets: no retained span contains a figure environment, a graphics-inclusion command, a PDF-page inclusion, a table or a TikZ picture; no image file is copied into this package, the package declares no asset dependency and no image file is redistributed with it. Licence: version arXiv:2607.28507v2 of this article is distributed under the Creative Commons Attribution 4.0 International licence, as recorded in the arXiv licence metadata for this exact version and shown on the abstract page https://arxiv.org/abs/2607.28507v2. A full rights notice is delivered at licenses/arxiv-2607.28507-CC-BY-4.0.txt. Characters: every retained excerpt is pure ASCII, so the delivered engine, XeLaTeX with its default Latin Modern text font, reports no missing character.
\begin{lemma}\label{lem:cuspnonzero}
 Suppose $S$ is nonempty. For $\chi\in \bigcap_{i\in S}\tau_{i,C}\chi_i\Lambda_{M_i},$ $\Ind_P^G \sigma_\chi\ge \sum_{i\in S} \Ind_{P_i}^G (\sigma_{i,\tau_{i,C}})_{\chi(\tau_{i,C}\chi_i)^{-1}}$ is a proper containment. In particular,
\begin{align*}
    \bigcap_{i\in S} \Ind_{K_{P_i}}^K V_{i,\tau_{i,C}}^\vee \neq 0.
\end{align*}
\end{lemma}

\begin{lemma}\label{lem:varying}
    There exists a family of holomorphic sections 
    \begin{align*}
        \mathfrak{H}=\{T_{S,C}:S\subseteq Y_{(M,\sigma)}, C\in \kappa_S\}\subseteq \left(\CC[\Lambda_M]\otimes \mathrm{End}_\CC (\Ind_{K_P}^K \sigma|_{K_M})\right)^{U_\sigma}
    \end{align*}
    such that for each $i\in Y_{(M,\sigma)},  v\in \mathrm{Ind}_{K_{P_i}}^K \sigma_{i,1}|_{K_{M_i}}$ and $\chi\in \chi_i\Lambda_{M_i}$
    \begin{align*}
        T_{S,C}(\chi)(v)=0
    \end{align*}
    for all $(S,C).$ 
    
    Moreover, if we let $\mathcal{H}$ denote the collection of all such families, the ideal in $\CC[\Lambda_M]$ generated by the set 
    \begin{align*}
        \{\Lambda_{M}\ni \chi \mapsto \langle T_{S,C}(\chi)(v),w \rangle: (v,w)\in \mathrm{Ind}_{K_{P}}^K \sigma|_{K_M}\times \mathrm{Ind}_{K_{P}}^K \sigma^\lor|_{K_M}, T_{S,C}\in \mathfrak{H}, \mathfrak{H}\in \mathcal{H}\}
    \end{align*}
    is unital. 
\end{lemma}

\begin{proof}      
  Choose any $U_\sigma$-eigenvector $T\in \mathrm{End}_\CC (\Ind_{K_P}^K \sigma|_{K_M})$ and $h\in \CC[\Lambda_M]^\times$ such that $T'_{\emptyset,\Lambda_M}:=hT\in (\CC[\Lambda_M]\otimes \mathrm{End}_\CC (\Ind_{K_P}^K \sigma|_{K_M}))^{U_\sigma}.$ If $J_{(M,\sigma)}$ is empty, take $$\mathfrak{H}:=\{T_{\emptyset,\Lambda_M}:=T'_{\emptyset,\Lambda_M}\},$$ and we are done.
 
 Thus we assume $J_{(M,\sigma)}$ is nonempty. Let $S\neq \emptyset$ and $C\in \kappa_S$. Then $C=U_\sigma E_C,$ where $E_C:=\bigcap_{i\in S} \tau_{i,C}\chi_i\Lambda_{M_i}.$  
Let $U'_\sigma\le U_\sigma$ be the subgroup that stabilizes $E_C$. Then $U'_\sigma$ is a subgroup of $\bigcap_{i\in S}\Lambda_{M_i}.$  Therefore, for all $\tau\in U'_{\sigma}$, we have identities of underlying vector spaces
\begin{align*}
    \bigcap_{i\in S}\mathrm{Ind}_{K_{P_i}}^{K}V_{i,\tau_{i,C}}^\lor=\bigcap_{i\in S}\mathrm{Ind}_{K_{P_i}}^{K}V_{i,\tau\tau_{i,C}}^\lor.
\end{align*}
By \hyperref[lem:cuspnonzero]{Lemma \ref{lem:cuspnonzero}} we can choose $0\neq w_{S,C}\in \bigcap_{i\in S}\mathrm{Ind}_{K_{P_i}}^{K}V_{i,\tau_{i,C}}^\lor \subseteq\mathrm{Ind}_{K_{P}}^{K} \sigma^\lor|_{K_M}$ so that
\begin{align*}
\langle v,w_{S,C}\rangle=0 \quad\textrm{ for any } v\in \sum_{i\in S}\mathrm{Ind}_{K_{P_i}}^{K} \sigma_{i,\tau_{i,C}}|_{K_{M_i}}.
\end{align*}
Choose $v_{S,C}\in \Ind_{K_P}^K\sigma|_{K_M}$ such that $\langle v_{S,C},w_{S,C}\rangle=1$. Consider the morphism $T'''_{S,C}\in \mathrm{End}_\CC (\Ind_{K_P}^K \sigma|_{K_M})$ corresponding to $v_{S,C}\otimes w_{S,C}\in \mathrm{Ind}_{K_P}^K\sigma|_{K_M}\otimes \mathrm{Ind}_{K_P}^K\sigma^\lor|_{K_M}.$ Choose a $U_\sigma'$-eigenvector 
\begin{align*}
    T\in \mathrm{span}_\CC\{ \tau.T'''_{S,C}: \tau\in U_{\sigma}'\}
\end{align*}
and $h\in \CC[\Lambda_M]^\times$ such that $T''_{S,C}:=hT\in (\CC[\Lambda_M]\otimes \mathrm{End}_\CC (\Ind_{K_P}^K \sigma|_{K_M}))^{U_\sigma'}.$ Then for any $\chi\in \Lambda_M$
\begin{align*}
    T_{S,C}''(\chi)(v)=0\quad\textrm{ for any } v\in \sum_{i\in S}\mathrm{Ind}_{K_{P_i}}^{K} \sigma_{i,\tau_{i,C}}|_{K_{M_i}},
\end{align*}
and there exists $(v_0,w_0)\in \Ind_{K_{P}}^K \sigma|_{K_{M}}\times \Ind_{K_{P}}^K \sigma^\lor|_{K_{M}}$ such that $\langle T_{S,C}''(\chi)(v_0),w_0\rangle\neq 0$ for any $\chi\in \Lambda_M.$

Let $\eta_{S,C}$ be a function in $\CC[\Lambda_M]^{U_\sigma'}$ whose vanishing locus $V(\eta_{S,C})$ in $\Lambda_M$ contains $U_\sigma \chi_i\Lambda_{M_i}-\tau_{i,C}\chi_i\Lambda_{M_i}$ for all $i\in S,$ and $V(\eta_{S,C})\cap E_C=\emptyset.$ Note that $V(\eta_{S,C})\supseteq C-E_C.$ Put
\begin{align*}
    T_{S,C}':=\sum_{\tau\in U_\sigma/U_{\sigma}'} \tau. \bigg(\eta_{S,C} T''_{S,C}\bigg).
\end{align*}
Then $T'_{S,C}\in \left(\CC[\Lambda_M]\otimes \mathrm{End}_\CC (\Ind_{K_P}^K \sigma|_{K_M})\right)^{U_\sigma}$ satisfies the following properties: 
\begin{itemize}
    \item For any $i\in S$ and $\chi\in \chi_i\Lambda_{M_i}$
\begin{align*}
    T_{S,C}'(\chi)(v)=0\quad\textrm{ for any } v\in \mathrm{Ind}_{K_{P_i}}^{K} \sigma_{i,1}|_{K_{M_i}}.
\end{align*}
    \item One has 
    \begin{align*}
        \langle T_{S,C}'(\chi)(v_0),w_0\rangle\neq 0.
    \end{align*}
    for any $\chi\in C$.
%     \item There exists $(v,w)\in \Ind_{K_{P}}^K \sigma|_{K_{M}}\times \Ind_{K_{P}}^K \sigma^\lor|_{K_{M}}$  such that 
% \begin{align*}
%      \langle T_{S,C}'(\chi)(v),w\rangle\neq 0.
% \end{align*}
\end{itemize}


 

For each $i\in Y_{(M,\sigma)},$ let $\{f_{i,1},\ldots,f_{i,p_i}\}\subseteq\CC[\Lambda_M]^{U_\sigma}$ be a generating set of the defining ideal of $U_\sigma\chi_i\Lambda_{M_i}$, viewed as a Zariski closed subspace of $U_\sigma\backslash \Lambda_M$. Fix a map $\xi:Y_{(M,\sigma)}\to \ZZ_{\geq 1}$ with $1\leq \xi(i)\leq p_i$. Define a partial order $\le$ on the finite set $\{(S,C): S\subseteq Y_{(M,\sigma)}, C\in \kappa_S\}$ by specifying $(S,C)\le (S',C')$ if and only if $C'\subseteq C$. Note that if $(S,C)\le (S',C'),$ then $C'\in \kappa_{S\cup S'}$ and $(S,C)\le (S\cup S', C').$ 
    
    
    If $(S,C)$ is not maximal (with respect to $\leq$) with $S\neq\emptyset$, let
    \begin{align*}
        h_{S,C,\xi}:=\prod_{\substack{i\in Y_{(M,\sigma)}\\ C\cap U_\sigma \chi_i\Lambda_{M_i}\subsetneq C}}f_{i,\xi(i)}=\prod_{\substack{i\notin S\\ C\cap U_\sigma \chi_i\Lambda_{M_i}\subsetneq C}} f_{i,\xi(i)}.
    \end{align*}
    For $(S,C)$ maximal, one has $C\cap U_\sigma\chi_i\Lambda_{M_i} = \emptyset$ for $i\not\in S$. We choose $h_{S,C,\xi}$ that is nonvanishing on $C$ and is zero on $U_\sigma\chi_i\Lambda_{M_i}$ for any $i\not\in S.$ 
    % Let $\mathcal{K}$ be the ideal in $\CC[\Lambda_M]^{U_\sigma}$ generated by 
    % \begin{align*}
    %     \{h_{S,C,\xi} :\emptyset\neq S\subseteq Y_{(M,\sigma)}, C\in \kappa_S,\,\xi\}.
    % \end{align*}
    % Note that the vanishing locus of $\mathcal{K}$ is empty. By Hilbert's Nullstellensatz, the ideal in $\CC[\Lambda_M]$ generated by $\mathcal{K}$ is unital. 
    
    Let $\{f_{1},\ldots,f_{p_\emptyset}\}\subseteq \CC[\Lambda_M]^{U_\sigma}$ be a generating set of the defining ideal of $\bigcup_{i\in Y_{(M,\sigma)}}U_\sigma \chi_i\Lambda_{M_i}.$ Extend the domain of $\xi$ to  $\{\emptyset\} \cup Y_{(M,\sigma)}$ with $1\le \xi(\emptyset)\le p_{\emptyset}.$ Let $h_{\emptyset,\Lambda_M,\xi}:=f_{\xi(\emptyset)}.$ Define for any $(S,C),\xi$
    \begin{align*}
       T_{S,C,\xi}:=h_{S,C,\xi}T'_{S,C}.
    \end{align*}
     Consider the families indexed by $\xi$
    \begin{align*}
        \mathfrak{H}_{\xi}:=\{ T_{S,C,\xi}: S\subseteq Y_{(M,\sigma)}, C\in \kappa_S\}.
    \end{align*}
    It follows by construction and Hilbert's Nullstellensatz that the ideal in $\CC[\Lambda_M]$ generated by the set 
    \begin{align*}
        \{\Lambda_{M}\ni \chi \mapsto \langle T_{S,C,\xi}(\chi)(v),w \rangle: (v,w)\in \mathrm{Ind}_{K_{P}}^K \sigma|_{K_M}\times \mathrm{Ind}_{K_{P}}^K \sigma^\lor|_{K_M}, T_{S,C,\xi}\in \mathfrak{H}_\xi, \xi\}
    \end{align*}
    is unital. In addition, given $(S,C),\xi,$ for any $i\in Y_{(M,\sigma)}, v\in \mathrm{Ind}_{K_{P_i}}^K \sigma_{i,1}|_{K_{M_i}}$ and $\chi\in \chi_i\Lambda_{M_i},$ we have $$T_{S,C}'(\chi)(v)=\delta_{i\notin S}T_{S,C}'(\chi)(v).$$
     On the other hand, by the definition of $h_{S,C,\xi}$ we have  $h_{S,C,\xi}(\chi)=0$ for $\chi\in \chi_i\Lambda_{M_i}$ if $i\not\in S$. Thus 
    \begin{align*}
    T_{S,C,\xi}(\chi)(v)=h_{S,C,\xi}(\chi)T_{S,C}'(\chi)(v)=0.
    \end{align*}
    This proves the lemma.
\end{proof}

\end{document}
